% !TEX program = xelatex
% 数学物理方法2 · §7.2 · 课件 p19 Assignments 1--4 · 四题中英答案（2026-11-09）
\documentclass[10pt]{ctexart}
%=====================================================================
% theme-violet.tex —— 讲义精讲 · 紫花主题
% 用法：\documentclass[11pt]{ctexart}
%       \input{assets/theme-violet}
%       \begin{document} ... \end{document}
% 只用 ctexart 自带能力 + xcolor/enumitem/fancyhdr/pifont/amsmath，
% 不依赖 tikz/tcolorbox，保证 MiKTeX 可直接编译。
%=====================================================================

\usepackage[a4paper,left=2.0cm,right=2.0cm,top=2.1cm,bottom=2.3cm,%
            headsep=12pt,footskip=22pt]{geometry}
\usepackage{amsmath,amssymb}
\usepackage{needspace}
\usepackage{xcolor}
\usepackage{array}
\usepackage{enumitem}
\usepackage{pifont}
\usepackage{fancyhdr}
\usepackage[hidelinks]{hyperref}

%---------------------------------------------------------------- 配色
\definecolor{violet}{RGB}{107,78,143}      % 主色：紫
\definecolor{violetdeep}{RGB}{63,44,92}    % 深紫：标题/强调
\definecolor{petal}{RGB}{183,159,208}      % 花瓣：装饰线、花饰
\definecolor{bandviolet}{RGB}{247,243,251} % 浅紫底：档案/提示框
\definecolor{bandcream}{RGB}{252,249,242}  % 米白底：易错/注意框
\definecolor{rulepurple}{RGB}{217,207,230} % 分隔线
\definecolor{inkgray}{RGB}{92,96,110}      % 次要文字

%---------------------------------------------------------------- 字体
% 有 Noto CJK 用 Noto，没有就退回 ctex 默认（Windows 下为宋体/黑体）
\IfFontExistsTF{Noto Serif CJK SC}{\setCJKmainfont{Noto Serif CJK SC}}{}
\IfFontExistsTF{Noto Sans CJK SC}{\setCJKsansfont{Noto Sans CJK SC}}{}

%---------------------------------------------------------------- 版式
\setlist{nosep,leftmargin=1.9em}
\linespread{1.28}
\setlength{\parskip}{3.5pt}
\setlength{\parindent}{0pt}
\setlength{\tabcolsep}{5pt}
\renewcommand{\arraystretch}{1.28}

%---------------------------------------------------------------- 花饰
% \ding{95} = 五瓣小花，\ding{96} = 星形花；都来自 pifont，必定存在
\newcommand{\orn}{\textcolor{petal}{\ding{95}}}          % 行内小标
\newcommand{\ornstar}{\textcolor{petal}{\ding{96}}}      % 星形花
\newcommand{\flowerrule}{%                                % 居中花饰分隔线
  \par\vspace{7pt}%
  \noindent\makebox[\linewidth]{\color{petal}%
    \rule{0.20\linewidth}{0.6pt}\hspace{8pt}{\small\ding{95}}\hspace{4pt}%
    \raisebox{1.1pt}{\tiny\ding{95}}\hspace{4pt}{\small\ding{95}}\hspace{8pt}%
    \rule{0.20\linewidth}{0.6pt}}%
  \par\vspace{7pt}}

%---------------------------------------------------------------- 标题与节头
\newcommand{\doctitle}[3]{%   #1 主标题  #2 副标题  #3 元信息（可用 \\ 分行）
  \begin{center}
    {\color{petal}\small \orn\enspace #2\enspace \ornstar}\\[7pt]
    {\LARGE\bfseries\color{violetdeep} #1}\\[10pt]
    {\color{petal}\rule{0.55\linewidth}{1pt}}\\[4pt]
    {\small\color{inkgray} #3}
  \end{center}\vspace{2pt}}

\newcommand{\hdr}[2]{%   #1 章节号（如 §0）  #2 标题
  \par\vspace{12pt}%
  \noindent\colorbox{violetdeep}{\color{white}\bfseries\small\;#1\;}\hspace{7pt}%
  {\color{violetdeep}\large\bfseries #2}\par\nobreak
  \vspace{3pt}{\color{petal}\hrule height 1.2pt}\vspace{6pt}}

\newcommand{\subhdr}[1]{%
  \par\vspace{6pt}%
  \noindent{\color{violet}\bfseries #1}\par\vspace{1pt}}

%---------------------------------------------------------------- 提示框
\newcommand{\callout}[1]{% 浅紫底：讲义档案、读法说明
  \par\medskip\noindent
  \fcolorbox{rulepurple}{bandviolet}{\parbox{\dimexpr\linewidth-2\fboxsep-2\fboxrule}{\vspace{2pt}#1\vspace{2pt}}}%
  \par\medskip}

\newcommand{\warnbox}[1]{% 米白底：易错、边界条件
  \par\medskip\noindent
  \fcolorbox{rulepurple}{bandcream}{\parbox{\dimexpr\linewidth-2\fboxsep-2\fboxrule}{\vspace{2pt}#1\vspace{2pt}}}%
  \par\medskip}

\newcommand{\gapitem}[2]{% #1 缺口类型（紫） #2 内容
  \par\vspace{2pt}\noindent
  \hangindent=1.4em\hangafter=1
  {\color{petal}\small\ding{96}}\enspace{\bfseries\color{violetdeep}#1}\enspace #2\par}

% 题干块（\qstem）：每道例题／作业的「详解」开头必须先复述题干与出处。
% 起因：2026-10-25 第五章讲义把七道作业的题面集中放在一页、详解另起一页，
% 读者翻到「第 4 题」时看不到条件，等于一道没有题干的题目。
% 规则：详解块内自带题干，读者不翻页就能读懂题在问什么。
\newcommand{\qstem}[1]{%
  \par\medskip\noindent
  \fcolorbox{violet}{bandviolet}{%
    \parbox{\dimexpr\linewidth-2\fboxsep-2\fboxrule}{%
      \vspace{2pt}%
      {\small\bfseries\color{violet}题干}\hspace{0.8em}#1%
      \vspace{2pt}}}%
  \par\medskip}

% v2: 保护题干与开头，长题可跨页，续页补上下文。
\newcommand{\qbreak}{\needspace{6\baselineskip}}
% 独立练习：ID、来源、题干、考点、答案标签（通常 ans:ID）。
\newcommand{\practice}[5]{%
  \qbreak\subhdr{练习 #1}\qstem{（#2）#3}%
  \textbf{考点}：#4\par
  详解见附录第~\pageref{#5}~页。\par\medskip}
% 附录开头：ID、来源、条件摘要；自动产生 ans:ID 标签。
\newcommand{\answerstart}[3]{%
  \qbreak\subhdr{#1\quad 详解}\label{ans:#1}\qstem{（#2）#3}}
\newcommand{\answercontinue}[2]{%
  \needspace{5\baselineskip}\subhdr{#1\quad 续解}%
  {\small\color{inkgray}本页接续：#2}\par\medskip}
% 延迟提供默认值，让旧讲义可在导言区保留自己的 \newcommand{\step}。
\AtBeginDocument{\providecommand{\step}[1]{\needspace{3\baselineskip}\par\medskip\textbf{#1}\par}}

% 前提旁注：只放在「需要理解的那一步」后面，给陈述 + 本题代入，不重证。
\newcommand{\aside}[1]{%
  \par\vspace{1pt}\noindent
  \hangindent=1.6em\hangafter=1
  {\small\color{inkgray}%
    \hspace*{0.6em}{\color{petal}$\triangleright$}\hspace{2pt}#1}%
  \par\vspace{2pt}}

%---------------------------------------------------------------- 表格
% 本机 MiKTeX 的 colortbl 与 array 版本不匹配（\rowcolor 会报错），
% 故表格统一用「上粗紫线 + 表头 + 中细线」的三线式，不依赖 colortbl。
\newcommand{\tabtop}{\par\vspace{3pt}{\color{violet}\hrule height 1.2pt}\vspace{1pt}}
\newcommand{\tabmid}{\par\vspace{1pt}{\color{petal}\hrule height 0.5pt}\vspace{1pt}}
\newcommand{\tabbot}{\par\vspace{3pt}{\color{rulepurple}\hrule height 0.5pt}\vspace{8pt}}
\newcommand{\thead}[1]{{\bfseries\color{violetdeep}#1}}

%---------------------------------------------------------------- 页眉页脚
\pagestyle{fancy}
\fancyhf{}
\renewcommand{\headrulewidth}{0.4pt}
\renewcommand{\headrule}{\hbox to\headwidth{\color{rulepurple}\leaders\hrule height \headrulewidth\hfill}}
\fancyhead[L]{\small\color{inkgray}\songti \doctag}
\fancyhead[R]{\small\color{inkgray}\songti \coursetag}
\fancyfoot[C]{\small\color{inkgray}\orn\enspace{\color{violet}\thepage}\enspace\orn}

% 每篇讲义在导言区设这两个值
\newcommand{\doctag}{}
\newcommand{\coursetag}{}
\newcommand{\settags}[2]{\renewcommand{\doctag}{#1}\renewcommand{\coursetag}{#2}}

%---------------------------------------------------------------- 数学记号
\newcommand{\vect}[1]{\boldsymbol{#1}}
\newcommand{\dd}{\,\mathrm{d}}
\newcommand{\ddx}{\frac{\mathrm{d}}{\mathrm{d}x}}
\addtolength{\jot}{2pt}



\settags{§7.2 作业 · 四题中英答案}{数学物理方法2}

\geometry{left=1.8cm,right=1.8cm,top=1.8cm,bottom=1.9cm,headsep=9pt,footskip=16pt}
\linespread{1.22}
\setlength{\headheight}{21pt}
\setlength{\parskip}{0pt}
\setlength{\abovedisplayskip}{4pt plus 1pt minus 1pt}
\setlength{\belowdisplayskip}{4pt plus 1pt minus 1pt}
\newcommand{\ok}{\textcolor{violet}{\ensuremath{\checkmark}}}
\newcommand{\en}[1]{\textit{#1}}
\newcommand{\tagline}[1]{{\small\color{inkgray}#1}}

\newcommand{\chdr}[2]{%
  \needspace{5\baselineskip}%
  \par\vspace{5pt}%
  \noindent\colorbox{violetdeep}{\color{white}\bfseries\small\;#1\;}\hspace{7pt}%
  {\color{violetdeep}\bfseries #2}\par\nobreak
  \vspace{1pt}{\color{petal}\hrule height 1pt}\vspace{4pt}}
\let\hdr\chdr

% 讲块（题目前）：为什么这么做、用哪条判据
\newcommand{\explain}[1]{%
  \par\medskip\noindent
  \fcolorbox{rulepurple}{bandcream}{\parbox{\dimexpr\linewidth-2\fboxsep-2\fboxrule}{%
    \vspace{2pt}{\color{violetdeep}\bfseries 讲解}\quad #1\vspace{2pt}}}%
  \par\medskip}

% 答案块（题目后）：中英并列
\newcommand{\answerbox}[2]{%
  \par\medskip\noindent
  \fcolorbox{violet}{bandviolet}{\parbox{\dimexpr\linewidth-2\fboxsep-2\fboxrule}{%
    \vspace{2pt}{\small\bfseries\color{violet}答案}\hspace{0.8em}#1\par\vspace{3pt}%
    {\small\bfseries\color{violet}Answer}\hspace{0.8em}\en{#2}\vspace{2pt}}}%
  \par\medskip}

\newcommand{\rodfig}[2]{%
  \begin{center}\small
  \begin{tabular}{c@{}c@{}c}
  #1 & {\setlength{\fboxsep}{0pt}\fcolorbox{petal}{petal!45}{\rule{0pt}{0.38cm}\hspace{2.6cm}}} & #2\\[1pt]
   & {\color{inkgray}\makebox[2.6cm][s]{$x=0$\hfill$x=l$}} &
  \end{tabular}
  \end{center}}

\begin{document}

\doctitle{§7.2 作业四题：定解条件怎么写}
{课件 p19 Assignments 1--4\ ·\ 每题按“讲解 $\to$ 题目 $\to$ 解答 $\to$ 答案（中英）”排列}
{依据课件 p19 题面、p14 力边界条件、p15 热流边界条件\quad 2026-11-09\quad 数学物理方法\,2}

\hdr{§0}{两条判据，四个符号}

四道题只用到两件事。第一件是\textbf{边界的类型}：给定边界面上的值（固定端、给定温度）是第一类；给定法向导数（受力、给定热流、绝热）是第二类。第二件是\textbf{外法向}：$x=0$ 处 $n$ 沿 $-x$，$x=l$ 处 $n$ 沿 $+x$，所以

\[
\left.\frac{\partial u}{\partial n}\right|_{x=0}=-\left.\frac{\partial u}{\partial x}\right|_{x=0},
\qquad
\left.\frac{\partial u}{\partial n}\right|_{x=l}=+\left.\frac{\partial u}{\partial x}\right|_{x=l}.
\]

把物理量翻译成法向导数，靠的是两条定律（课件 p14、p15）：

\[
\text{杆受力：}\ \frac{f_n}{S}=E\frac{\partial u}{\partial n},
\qquad
\text{热流：}\ q_n=-k\frac{\partial u}{\partial n},
\]

$f_n$ 是外力沿\textbf{外法向}的分量（向外为正），$q_n$ 是热流沿外法向的分量（流出为正）。注意第二式自带一个负号，第四式的“流入”与第三式的“受压”箭头同向、符号却相反，根源就在这个负号上。

\hdr{题 1}{左端固定，右端受外向拉力}

\explain{先看两端各“给”了什么。左端写“fixed”，给的是位移被按死，属于第一类，条件是函数值等于零；右端写“a normal force acts outward”，给的是力，属于第二类，条件是法向导数等于力除以刚度。右端的外法向朝 $+x$，力又朝外，所以 $f_n=+F(t)$，代数上不用再翻号——这是四题里最顺的一道。}

\qstem{\en{The left end of a rod (whose length is $l$) is fixed, and a normal force $F(t)=F_0\sin\omega t$ acts outward on the right end of the rod.}\\[3pt]
\textbf{中}（p19 第 1 题）：长为 $l$ 的杆，左端固定，右端受一外向法向力 $F(t)=F_0\sin\omega t$。求边界条件。}

\subhdr{解答过程}

\begin{enumerate}
\item 左端 $x=0$：固定端，位移恒为零，得第一类齐次条件 $u(0,t)=0$。
\item 右端 $x=l$：受外力。外法向 $n$ 指向 $+x$，力沿 $+x$ 向外，故 $f_n=+F_0\sin\omega t$。
\item 代入 Hooke 定律 $\dfrac{f_n}{S}=E\dfrac{\partial u}{\partial n}$，并用右端 $\partial u/\partial n=\partial u/\partial x$：
\[
\frac{F_0\sin\omega t}{S}=E\left.\frac{\partial u}{\partial x}\right|_{x=l}.
\]
\item 两边乘 $S$，写成作业常用的形式 $ES\,u_x(l,t)=F_0\sin\omega t$。
\end{enumerate}

\answerbox{$u(0,t)=0$（左端固定）；$\;ES\,u_x(l,t)=F_0\sin\omega t$，其中 $E$ 为杨氏模量、$S$ 为横截面积。（也可写成 $u_x|_{x=l}=F_0\sin\omega t/(ES)$。）}
{Boundary conditions: $u(0,t)=0$ (the left end is fixed) and $ES\,u_x(l,t)=F_0\sin\omega t$ (an outward normal force acts on the right end), where $E$ is Young's modulus and $S$ the cross-sectional area.}

\warnbox{\textbf{自检}：令 $F_0=0$，右端变成 $\left.u_x\right|_{x=l}=0$，正是 p14 的 “Free!”（自由端）\ok{}。反例：写成 $u(l,t)=0$ 就把右端当成了固定端，与“受拉力”矛盾。}

\hdr{题 2}{两端固定的弦，$t=0$ 时一点受外力}

\explain{这题问的是初始条件，不是边界条件——两端固定只是背景。“波动方程对时间二阶，所以初值要两个”：初始位移、初始速度。弦原本静止，位移初值为零；t=0 的一点受力，本质是一个\textbf{冲量}，它把速度在某一点上“打”成一个 $\delta$。判断用冲量—动量定理，量纲会替你检验：只有把 $F_0$ 当冲量（N$\cdot$s），$(F_0/\rho)\delta$ 才是速度。}

\qstem{\en{A string, whose length is $l$, is fixed at its end points. At $t=0$, an external force $F_0$ acts on the point $x=h$.}\\[3pt]
\textbf{中}（p19 第 2 题）：长为 $l$ 的弦，两端固定；$t=0$ 时在 $x=h$ 处受一外力 $F_0$。求初始条件。}

\subhdr{解答过程}

\begin{enumerate}
\item 先把问题分开：两端固定属于边界条件（$u(0,t)=u(l,t)=0$），本题不问；要写的是 $t=0$ 的两个初始条件。
\item 初始位移：$t<0$ 时弦静止在平衡位置，故 $u(x,0)=0$。
\item 初始速度：取 $x=h$ 附近长 $\Delta x$ 的一小段，冲量 $F_0$ 使它获得动量 $\rho\,\Delta x\,u_t$，于是 $u_t=\dfrac{F_0}{\rho\,\Delta x}$；令 $\Delta x\to0$，该速度集中在一点，写成 $\delta$ 函数：
\[
u_t(x,0)=\frac{F_0}{\rho}\,\delta(x-h).
\]
\item 量纲核对：$\dfrac{F_0}{\rho}$ 的单位是 $\dfrac{\mathrm{N\cdot s}}{\mathrm{kg/m}}=\mathrm{m^2/s}$，再乘 $\delta$ 的 $\mathrm{m^{-1}}$ 得 $\mathrm{m/s}$，正是速度 \ok{}。
\end{enumerate}

\answerbox{$u(x,0)=0$；$\;u_t(x,0)=\dfrac{F_0}{\rho}\,\delta(x-h)$，其中 $\rho$ 为弦的线密度，$F_0$ 按冲量理解。}
{Initial conditions: $u(x,0)=0$ and $u_t(x,0)=\dfrac{F_0}{\rho}\delta(x-h)$, where $\rho$ is the linear density of the string and $F_0$ is the impulse.}

\warnbox{\textbf{一个必须说清的约定}：若把 $F_0$ 当成普通力（单位 N）而不是冲量，它就不属于初始条件，而属于方程右端的源项 $\dfrac{F_0}{\rho}\delta(x-h)$，此时初条件只剩 $u(x,0)=0$、$u_t(x,0)=0$。两种理解只差“力持续多久”，量纲一眼可分。}

\hdr{题 3}{两端受压的杆}

\explain{两端给的都是力，所以两端都是第二类；真正要小心的是方向。图里两个箭头都指向杆内，是压力。外法向在两端相反，“向内”在两端对应的空间方向也相反，于是最省事的写法是先统一写成 $f_n$（沿外法向的分量）：压力与外法向反向，$f_n=-F_0$，两端一样；最后换到 $\partial u/\partial x$ 时，左端才翻号。}

\qstem{\en{A normal force $F_0$ acts on the ends of the rod (whose length is $l$).}\\[3pt]
\textbf{中}（p19 第 3 题）：长为 $l$ 的杆，两端各受一法向力 $F_0$（图中箭头指向杆内，即两端受压）。求边界条件。}

\rodfig{$F_0\;\longrightarrow$}{$\longleftarrow\;F_0$}

\subhdr{解答过程}

\begin{enumerate}
\item 两端都只给力、不给位移，故两端都是第二类边界条件。
\item 取外法向：左端 $n$ 沿 $-x$，右端 $n$ 沿 $+x$。压力方向与 $n$ 相反，所以两端都有 $f_n=-F_0$。
\item 代入 $\dfrac{f_n}{S}=E\dfrac{\partial u}{\partial n}$，得两端统一的形式
\[
\left.\frac{\partial u}{\partial n}\right|_{\text{两端}}=-\frac{F_0}{ES}.
\]
\item 换成 $x$ 导数：左端 $\partial u/\partial n=-\partial u/\partial x$，故 $-\left.u_x\right|_{x=0}=-\dfrac{F_0}{ES}$，即 $\left.u_x\right|_{x=0}=+\dfrac{F_0}{ES}$；右端 $\partial u/\partial n=+\partial u/\partial x$，直接得 $\left.u_x\right|_{x=l}=-\dfrac{F_0}{ES}$。
\end{enumerate}

\answerbox{$\left.u_x\right|_{x=0}=+\dfrac{F_0}{ES}$，$\;\left.u_x\right|_{x=l}=-\dfrac{F_0}{ES}$；等价地，外法向形式 $\left.\dfrac{\partial u}{\partial n}\right|_{\text{两端}}=-\dfrac{F_0}{ES}$。}
{$\left.u_x\right|_{x=0}=+F_0/(ES)$ and $\left.u_x\right|_{x=l}=-F_0/(ES)$; equivalently, in terms of the outward normal, $\partial u/\partial n=-F_0/(ES)$ at both ends.}

\warnbox{\textbf{与 p14 对照}：p14 是“外力向外”，得左端 $-\frac{f}{ES}$、右端 $+\frac{f}{ES}$；本题力向内，整组变号。\textbf{不要}只翻一端——若只把右端改成负号、左端照抄 p14，就得到两端不同向的怪结果。}

\hdr{题 4}{两端流入热流的杆}

\explain{给的是热流密度，第二类；判据是 Fourier 定律，它自带负号。箭头朝内是“流入”，流入意味着热流逆着外法向，所以 $q_n=-q_0$，两端一样。换成 $\partial u/\partial x$ 后是“左负右正”，与题 3 正好相反——不是算错了，而是 Hooke 定律没有负号、Fourier 定律有负号。}

\qstem{\en{The heat flux into the ends of the rod (whose length is $l$) is $q_0$.}\\[3pt]
\textbf{中}（p19 第 4 题）：长为 $l$ 的杆，两端各流入热流密度 $q_0$。求边界条件。}

\rodfig{$q_0\;\longrightarrow$}{$\longleftarrow\;q_0$}

\subhdr{解答过程}

\begin{enumerate}
\item 未知量 $u$ 是温度，给定的是端面热流密度，故为第二类边界条件。
\item 取外法向 $n$，Fourier 定律 $q_n=-k\dfrac{\partial u}{\partial n}$；热流流入杆内即逆着 $n$，故 $q_n=-q_0$，两端相同。
\item 于是两端统一地有 $\dfrac{\partial u}{\partial n}=\dfrac{q_0}{k}$。
\item 换成 $x$ 导数：左端 $-\left.u_x\right|_{x=0}=\dfrac{q_0}{k}$，故 $\left.u_x\right|_{x=0}=-\dfrac{q_0}{k}$；右端 $\left.u_x\right|_{x=l}=\dfrac{q_0}{k}$ 直接得出。
\end{enumerate}

\answerbox{$\left.u_x\right|_{x=0}=-\dfrac{q_0}{k}$，$\;\left.u_x\right|_{x=l}=+\dfrac{q_0}{k}$，$k$ 为热导率；等价地 $-k\left.u_x\right|_{x=0}=q_0$、$-k\left.u_x\right|_{x=l}=-q_0$（后一式读作：$+x$ 方向净热流为 $-q_0$，真实热流沿 $-x$ 流入）。}
{$\left.u_x\right|_{x=0}=-q_0/k$ and $\left.u_x\right|_{x=l}=+q_0/k$ with $k$ the thermal conductivity; equivalently the inward heat flux at both ends is $q_0$.}

\warnbox{\textbf{自检}：两端吸热，端部比内部热，温度从左端往里降、从右端往里降，故左端 $u_x<0$、右端 $u_x>0$ \ok{}；$q_0=0$ 时两端绝热（$\left.u_x\right|=0$），回到 p15 的 adiabatic \ok{}。}

\hdr{§5}{四题答案速查（中英对照）}

\par\smallskip
\begin{itemize}
\item \textbf{题 1}（左端固定，右端外向力 $F_0\sin\omega t$）：$u(0,t)=0$；$ES\,u_x(l,t)=F_0\sin\omega t$。\en{EN: $u(0,t)=0$; $ES\,u_x(l,t)=F_0\sin\omega t$.}
\item \textbf{题 2}（两端固定弦，$t=0$ 于 $x=h$ 受冲量 $F_0$）：$u(x,0)=0$；$u_t(x,0)=\frac{F_0}{\rho}\delta(x-h)$。\en{EN: $u(x,0)=0$; $u_t(x,0)=\frac{F_0}{\rho}\delta(x-h)$.}
\item \textbf{题 3}（两端受压 $F_0$）：$u_x(0)=\frac{F_0}{ES}$，$u_x(l)=-\frac{F_0}{ES}$。\en{EN: $u_x(0)=F_0/(ES)$; $u_x(l)=-F_0/(ES)$.}
\item \textbf{题 4}（两端流入热流 $q_0$）：$u_x(0)=-\frac{q_0}{k}$，$u_x(l)=+\frac{q_0}{k}$。\en{EN: $u_x(0)=-q_0/k$; $u_x(l)=+q_0/k$.}
\end{itemize}

\hdr{§6}{四个错法，对号入座}

\gapitem{题 1 把右端当固定端}{写成 $u(l,t)=0$ 就变成“两端固定”。题目给的是力，力只能进第二类条件（法向导数），不能变成位移为零。}

\gapitem{题 2 只写位移、漏掉速度}{波动方程对 $t$ 是二阶的，一个初始条件不够。漏掉 $u_t(x,0)$ 等于只给了形状、没给初速度，解不唯一。}

\gapitem{题 3 只翻一端的号}{“两端受压”意味着两端都反向。只把右端改成负号、左端照抄 p14 的 $-\frac{f}{ES}$，结果两端同号，与“两端受压”的对称性矛盾。}

\gapitem{题 4 忘记 $k$}{写成 $u_x=\mp q_0$ 量纲不成立：$q_0$ 是 $\mathrm{W/m^2}$，$u_x$ 是 $\mathrm{K/m}$，中间必须有热导率 $k$ 才能换算。}

\hdr{§7}{符号与量纲速查}

{\small
\tabtop
\begin{tabular}{@{}c p{4.6cm} p{5.0cm} p{4.4cm}@{}}
\thead{符号} & \thead{含义} & \thead{单位} & \thead{出现} \\
\hline
$u$ & 位移（题 1--3）／温度（题 4） & $\mathrm{m}$ ／ $\mathrm{K}$ & 四题 \\
$E$ & 杨氏模量 & $\mathrm{Pa}=\mathrm{N/m^2}$ & 题 1、3 \\
$S$ & 杆的横截面积 & $\mathrm{m^2}$ & 题 1、3 \\
$f_n$ & 外力沿外法向的分量（向外为正） & $\mathrm{N}$ & 题 1、3 \\
$F_0$ & 外力幅值；题 2 中按冲量理解 & $\mathrm{N}$ ／ $\mathrm{N\cdot s}$ & 题 1--3 \\
$q_0$ & 热流密度 & $\mathrm{W/m^2}$ & 题 4 \\
$k$ & 热导率 & $\mathrm{W/(m\cdot K)}$ & 题 4 \\
$\rho$ & 弦的线密度 & $\mathrm{kg/m}$ & 题 2 \\
$\delta(x-h)$ & 集中在一点的点源 & $\mathrm{m^{-1}}$ & 题 2 \\
\end{tabular}
\tabbot
}

\callout{\textbf{一句话带走}：四道题只有两种动作——\emph{给值的写第一类，给力或热流的写第二类}；写第二类时先把物理量写成“沿外法向的分量”，再用 Hooke（不带负号）或 Fourier（带负号）换成 $\partial u/\partial n$，最后按端面把 $\partial u/\partial n$ 换回 $\partial u/\partial x$。}

\vspace{2pt}
\tagline{第 4 题的展开版（三条等价写法、稳态讨论、与 p14 的同构对照）另见本目录《作业 4 解答卡（20261109）》。}

\end{document}
